% !PRINTED 636
\weberpage{636}
%<<CONT
then by \S.\,172, III. the decomposition of $p$ into prime factors in the field $\Omega_m$ is the same as in the field $\Omega_{m_1}$, and we have
\[
\mathfrak{p}_n = \mathfrak{p}_{n+sm_1},
\]
if $s$ is an arbitrary integer. We now set
\weq{18}{t = n + s m_1}
and let $n$ run through the sequence of the positive integers $<m_1$ not divisible by $q$, and $s$ through the sequence of numbers $0,1,2,\ldots,m_2-1$; then $t$ runs through a complete system of residues of $m$ not divisible by $q$, and indeed the completely determinate system whose elements are all $<m$. We then have by \S.\,172
\weq{19}{p = \prod^{n} \mathfrak{p}_n.}

If, further, $t'$ is a number which satisfies the congruence
\weq{20}{t\,t' \equiv 1\ (\mathrm{mod}\ m),}
then, if $n'$ is determined by the condition $n n' \equiv 1\ (\mathrm{mod}\ m_1)$,
\weq{21}{t' \equiv n'\ (\mathrm{mod}\ m_1).}

According to this the product
\[
\prod^{t} \mathfrak{p}_{t'}^{\,t}
\]
can be determined by the formula (16), which, applied to the field $\Omega_{m_1}$, consisting of all rational functions of the $m_1$th root of unity $r_1 = r^{m_2}$, yields the formula:
\weq{22}{(r_1,\eta)^{m_1} = \prod^{n} \mathfrak{p}_{n'}^{\,n}.}

But by (18) and (21)
\[
\prod^{t} \mathfrak{p}_{t'}^{\,t} = \prod^{n} \mathfrak{p}_{n'}^{\,n m_2} \prod^{n'}\prod^{s} \mathfrak{p}_{n'}^{\,s m_1},
\]
and since $n'$, according to the modulus $m_1$, runs through the same sequence of numbers as $n$, we have by (19):
\[
\prod^{s}\prod^{n'} \mathfrak{p}_{n'}^{\,s m_1} = p^{m_1 \Sigma s} = p^{\frac{1}{2}m(m_2-1)},
\]
and there results:
\[
\prod^{t} \mathfrak{p}_{t'}^{\,t} = \left[p^{\frac{1}{2}(m_2-1)}\,(r^{m_2},\eta)\right]^m.
\]

If we set
\weq{23}{\varrho = p^{\frac{1}{2}(m_2-1)}\,(r_1,\eta), \qquad \varrho_n = p^{\frac{1}{2}(m_2-1)}\,(r_1^{\,n},\eta),}
then $\varrho$ is a \gesperrt{cyclotomic number}, since indeed $\sqrt{p}$ too, as a value
%CONT>>
